% ============================================================
%  H. Høffding, "Vort Hjem. En Indledning"
%  Indledning til: Fru Emma Gad (red.),
%  Vort Hjem, Bind I.
%  København: Det Nordiske Forlag, Ernst Bojesen, 1903.
%  TRANSCRIPTION (Danish)
% ============================================================
\documentclass[12pt,a4paper]{article}
\usepackage[utf8]{inputenc}
\usepackage[T1]{fontenc}
\usepackage{libertinus}
\usepackage{libertinust1math}
\usepackage[danish]{babel}
\usepackage[protrusion=true,expansion=true]{microtype}
\usepackage{setspace}
\usepackage[margin=1.2in]{geometry}
\usepackage{fancyhdr}
\usepackage{hyperref}

\hypersetup{
  pdftitle={Vort Hjem. En Indledning},
  pdfauthor={H. Høffding},
  hidelinks
}

\pagestyle{fancy}
\fancyhf{}
\rhead{Høffding, \emph{Vort Hjem} (1903)}
\cfoot{\thepage}

\setstretch{1.3}

\title{\textbf{Vort Hjem}\\[6pt]
  \large En Indledning}
\author{Harald Høffding}
\date{Indledning til: Fru Emma Gad (red.),
  \emph{Vort Hjem}, Bind I\\
  København: Det Nordiske Forlag, Ernst Bojesen, 1903}

\begin{document}
\maketitle
\thispagestyle{fancy}

\bigskip

\section*{I.}

\noindent Ordet Hjem vækker Forestillinger om Inderlighed, Fylde og Hvile.
Om Inderlighed --- thi skønt Hjemmet ude fra set er et Sted som alle
andre Steder, er det dog for os det Sted, hvor vi modtog alt, hvad
Skæbnen bragte os, den Vig, til hvilken Virkningerne fra den ydre
Verdens store Hav maatte forplante sig, for at vi skulde mærke og
værdisætte dem; og det er tillige det Sted, fra hvilket de Virkninger,
vi selv maatte have øvet, ere udgaaede. I al vor Virken og Liden,
al vor Given og Modtagen væver derfor Hjemmets Stemning sig ind
og danner en Baggrund, som vi maaske først blive os klart bevidst,
naar den trues af Opløsning eller dog undergaar dybere gaaende
Ændring. Og hermed hænger ogsaa Fylden sammen. Thi det er en
stor Kreds af Erindringer og Vaner, egentligt hele vort Livs Indhold,
der hænger sammen med denne centrale Forestilling. Hjemmet er
som Nettet, der ikke kan drages op, uden at et Mylr af levende
Indhold følger med. Hvilen endelig beror baade paa Inderlig\-heden
og Fylden. Hvile kan ikke findes i det fremmede og nye, kun i
det, der er kendt og prøvet, og hvis Tilegnelse ingen Anspændelse
kræver. Det hjemlige er lempet efter vort Væsens Former; alle
Folder og Fuger i vor Natur finde der et Leje, som passer til dem.
Naar det hjemlige atter dukker op for os, er det, som om det aldrig
havde været borte; let og hurtigt føje vi os ind i de forlængst til-
lempede Rammer. Og den Rigdom af Vaner og Erindringer, Hjemmet
bærer, bringer ogsaa Hvile. Paa en skarp Od kan man ikke hvile;
der maa være en Mangfoldighed i Indholdet, for at den stadige
Dvælen ved et og det samme ikke skal trætte og sløve.

Et Hjem har man, saa snart man har noget at hælde sit Hoved
til. Hjemmet forudsætter Fasthed, Stadighed, Gentagelse. Det
hjemlige staar i Modsætning til det nye, fremmede og skiftende. Naar
Ræven har sin Hule og Fuglen sin Rede, have de et Hjem. De have
noget at vende tilbage til, noget, der venter paa dem, noget, der kan
genkendes, og som staar fast under Forandringernes Bølgegang. Et
Hjem have vi derfor ogsaa i vort eget Indre, naar vor Bevidsteds
Indhold ikke uafbrudt skifter, men vi have Grundtanker og Grund-
stemninger, som udgøre vort aandelige Stade og betinge den Virk-
ning, de nye Erfaringer gøre paa os, ja ere Forudsætningen for, at
noget kan komme til at staa som nyt for os. I disse Grundtankers
og Grundstemningers Kreds maa alt optages, hvad vi virkeligt skulle
kalde vort aandelige Eje.

Det er mere end en vilkaarlig Sammenligning, vi her anstille,
naar vi stille det ydre, materielle Hjem og det indre, aandelige Hjem
sammen. Thi de bero begge paa de samme fysiologiske og psyko-
logiske Betingelser. Livet kan ikke bestaa uden en Rytme i de
indre og ydre Forhold. Optagen og Udskillen, Indaanding og Ud-
aanding, Hvile og Bevægelse maa afløse hinanden, for at Livet kan
bestaa. Altsaa Tilbagevenden til samme Tilstande er en organisk
Livsbetingelse. Og sjæleligt kræver Livets Gang faste Udgangs- og
Orienteringspunkter, stadige Baggrundstilstande, i Modsætning til
hvilke Oplevelserne kunne fremtræde i deres Ejendommelig\-hed
og dog saaledes, at en fælles Maalestok bliver mulig.

\section*{II.}

\noindent I den menneskelige Verden høre Hjem og Familie nøje sammen.
Dog kan der være et Hjem, hvor der ingen Familie er. Eneboeren
kan have sit Hjem, og paa de laveste Udviklingstrin af Menneskeliv,
vi kende, har Familien endnu ikke udsondret sig fra Horden eller
Klanen som et særeget Samfund. Den omvankende Horde vender
med Forkærlighed tilbage til de Steder, der tilbyde Næring og Tryg-
hed, og den fører i hvert Tilfælde sine Telte og sine Brugsgenstande
med sig omkring som Bestanddele af et Vandrehjem. Allerede her
betyder Hjemmet Inderlighed, Fylde og Hvile. Kampen for Livet,
som er en Kamp mod det fremmede, stanser i Teltet og paa de
kendte Græsgange og Jagtgrunde. Og Erindringerne spille allerede
her en Rolle, naar man færdes i samme Egne, hvor Forfædrene
færdedes. Om man ellers ingen faste Hjemsteder har, vil dog Fæ-
drenes Grave være et Centrum, man kredser om, og hvor man tager
Kampen op, selv om man ellers ingen Modstand gør mod frem-
trængende fremmede Skarer. Allerede her undergaar da Hjemmet
den betydningsfulde Metamorfose fra at være et materielt Nødvendigt
til at blive et ideelt Værdifuldt. Hvad der fra først af kun var et
fysisk Tilhold, kan siden blive det eneste Sted, hvor Livet synes
værd at leve, fordi det kun dér kan leves i Sammenhæng med For-
tidens Liv og som Fortsættelse af dette.

Af største Betydning i Hjemmets Historie blev Familiens Udvikling
til en særegen Art af Samfund. Familiens Kærne er Forbindelsen af
Moder og Barn. De to holde sammen og have deres Hjem hos hin-
anden, hvorledes det saa gaar med Familiens andre Elementer. Her
virker Naturforbindelsen, ikke blot den Forbindelse, der betinges ved,
at Trangen til Næring og Tryghed tilfredsstilles paa de samme
Steder. Fødselen er en Forlængelse af den organiske Vækst ud over
den individuelle Organisme. Hjemmet vokser altsaa bogstaveligt frem.
Rundt om paa Jorden finder man paa et vist Trin af Udvikling den
Ordning, at Mand og Hustru vedblive at tilhøre hver sin Stamme,
og at Børnene tilhøre Hustruens, ikke Mandens Stamme. Hvad der
har skabt Familien, saaledes som vi nu kende den, som et Samfund
af Mand, Kvinde og Børn, har dels været Mandens Magtlyst og Selv-
følelse, dels den faste Bolig. Manden betragtede Kvinden og de Børn,
hun havde født ham, som sin Ejendom; han ønskede at udnytte den
Arbejdskraft, de besad, og han følte Trang til at se sin Slægts Liv
fortsættes i sit Afkom. Den faste Bolig og det regelmæssige Samliv,
den medførte, afgav et Grundlag for Sæd og Skik, for Overlevering
og Erindring, og der blev nu rigeligere Lejlighed for ømmere Følel-
sers Opstaaen. Samlivet mildnede, underbyggede og humaniserede
den dyriske Trang og den raa Magtlyst. Den fælles Arne blev hellig,
fordi Forfædrenes Aander stadigt besøgte den. Og i denne Fred-
hellighed fik ogsaa den fremmede Del --- selv om han var en
Fjende --- en Helligdom at ty til. Der var skabt en Helligdom,
hvis Aand kunde brede sig til større Kredse.

Familien er paa laveste Trin kun et Middel for Livets Bestaaen
og Fortsættelse. Men ogsaa her danner der sig Livsinteresser og
Motiver af anden Art end de, der fra først af grundlagde Forbindel-
sen. Det er da ikke blot Næring, Tryghed og Slægtens Fortsættelse,
der søges; den hjemlige Stemning, som det faste Tilhold og de fælles
Overleveringer afføde, bliver i sig selv værdifuld, staar som et umiddel-
bart Gode, i hvilket Mennesket ønsker at leve. Fælles Skæbne og
fælles Arbejde befæster stadigt Følelsen af Sammenhøren. Selv om
Naturbaandet kun i Forholdet mellem Moder og Barn træder aldeles
øjensynligt frem, væves der under det fælles Livs Gang Baand mellem
alle Familiens Medlemmer, der ofte i Styrke kunne tage det op med
den blot fysiske Forbindelse.

Selv efter at Familien har udsondret sig fra Horden eller Klanen
og er bleven et Samfund for sig, baaret af ejendommelige Kræfter,
træder den endnu ikke i den Modsætning til andre Arter af Samfund,
som vi kende fra det moderne Liv. Den har flere Formaal end den
moderne Familie, idet den tillige er en lille Stat og en lille Kirke,
Retssamfund og Kultursamfund paa én Gang. Fædre- og Forældre-
magten raader som højeste Myndighed, og selv om en Stat eksisterer,
danner Familiens Tærskel en Grænse for dens Indflydelse. Religionen,
der paa dette Trin er Indbegrebet af al Kultur, er fra først af Fa-
miliens Sag: i Familien dyrkes de afdødes Aander fra Slægt til Slægt;
til deres Ære holdes Arneilden vedlige, og deres Bud er den moralske
Lovs Indhold. Først under den fremskridende Udvikling ophører
Familien at have saa omfattende Formaal; Staten inddrager alle For-
hold mellem Individer under sin Magt og sit Opsyn, og der dannes
særegne Samfund til Pleje af materiel og ideel Kultur. Tilbage
bliver da for Familien Plejen af de første Spirer til nyt menneskeligt
Liv og af de inderlige Følelser, der kunne knytte Mennesker sammen
afset fra alle Kulturbestræbelser og alle Retsforhold. Herved forringes
Familiens Betydning ingenlunde. Tværtimod staar Familieforholdet
ved sin Uafhængighed af alle de Forskelle og Modsætninger, der i
større Kredse dele Menneskene, og ved sin Sammenhæng med de
Naturinstinkter, der sikre Fortsættelsen af Slægtens Liv, som det
stærkeste og inderligste Samfund, der afgiver en Maalestok for alle
andre Samfunds Udvikling; jo inderligere og stærkere andre sociale
Forbindelser blive, des mere betegnes de ved Udtryk, der ere hentede
fra Familiens. I Familien kan den enkelte leve efter sin hele Per-
sonlighed, medens Kultursamfundet og Staten kun har Brug for en-
kelte Sider af hans Væsen. Og Forbindelsen bunder her i det
ubevidste Naturliv, medens andre Samfund grunde sig paa halvt
eller helt bevidste Formaal og Interesser.

Selv hvor Arbejdsdelingen fører det med sig, at der indtræder en
Modsætning mellem Familien paa den ene Side, Kultursamfund og
Stat paa den anden Side, bevarer Familien derfor sin Betydning,
eller den faar da først ret sin Betydning. Den bliver nu et Central-
sted, hvor det, der bevirkes ved Arbejdet i Kultursamfund og Stat,
kan værdsættes og nydes. Og den bliver et Arnested, hvor de
Kræfter, der skulle virke i Kultursamfund og Stat, kunne spire og
vokse. Den repræsenterer Inderlighed, Samling, Tryghed og
Hvile, til hvilke der atter og atter kan vendes tilbage fra den udad-
vendte, delte, kæmpende og anspændte Virken i de større Kredse.

\section*{III.}

\noindent Vi have skildret Hjemmet efter dets Ideal saavel som efter dets
historiske Udvikling. Men det viste sig, at saare mange forskellige
Kræfter og Betingelser forudsættes, for at Idealet her kan naas, og det
er da at vente, at Idealet ikke altid naas. Og hvis et gammelt Ord har
Ret, at jo bedre noget er, des slettere bliver det, naar det forvanskes,
saa maa vi vente, at Hjemmet under uheldige Forhold vil kue Livet
lige saa meget, som det under lykkelige Forhold kan fremme det.

Lad os først noget nøjere undersøge, hvorpaa Hjemmets blivende
Betydning og Berettigelse grunder sig. Begrundelsen ligger i en lig-
nende Tankegang som den, der begrunder den enkelte Personligheds
Ret og Frihed. For Samfundets egen Skyld maa den enkeltes størst
mulige Selvstændighed og Frihed sikres. Thi i den enkeltes Indre
ligger Spiren til alt nyt og stort. Jo kraftigere Livet rører sig der,
des kraftigere vil det kunne røre sig i de store Kredse. De enkelte
Personligheder ere Kildestederne for Samfundets Liv, lige som de ere
de Centralsteder, til hvilke alle af Samfundslivet udspringende Virk-
ninger forplante sig. Dér føles Livets Værdi, og derfra udgaa de
Bestræbelser, der kunne forhøje denne Værdi. Paa lignende Maade
forholde de mindre og snævrere Samfund sig til de større. Livet føres
og føles i hine stærkere og inderligere end i disse. Hvad der ikke
paaagtes, og hvad der ikke skærmes i de større Kredse, finder Op-
mærksomhed og Omsorg i de mindre. Derfor kræver de store Sam-
funds --- Statens og Kultursamfundets --- egen Interesse, at de mindre
hævdes i deres Ejendommelighed og Selvstændighed. Paa det om-
vendte Forhold mellem Livets Styrke og Omfang grunder sig saaledes
ogsaa Hjemmets Betydning. Allerede Aristoteles har i sin Kritik af
Platon's »Stat«, hvor Familien for de højere Klassers Vedkommende
skulde ophæves til Fordel for et almindeligt Slægtskab mellem alle
Individer, støttet sig til denne Lov. Det vil, siger han, blive en
vandet Kærlighed, som saaledes efter Ophævelsen af de snævrere og
inderligere Forhold skulde kunne forene alle; --- den vil næppe være
til at mærke, lige saa lidt som Iblanding af en lille Smule Sukker i
en stor Mængde Vand gør Indtryk paa Smagen. Det er fuldstændigt
rigtigt, siger Aristoteles videre, at Staten skal være en Enhed; men
udsletter man al Forskel og Mangfoldighed, hører Staten op at være,
lige som Symfonien hører op, naar man kun lader én Stemme lyde,
og Rytmen, naar man kun har én eneste Versefod. --- Den store
græske Statsfilosof har hermed klart vist Betydningen af de smaa
Samfunds Bestaaen inden for det store.

Men for at forstaa de Tvivl, der kunne opstaa om Hjemmets
Værdi, maa vi mærke os, at Loven om det omvendte Forhold mel-
lem Styrke og Omfang ikke er den eneste Lov for menneskelige
Interessers Udvikling. Der viser sig nemlig paa Følelseslivets Om-
raade ogsaa et omvendt Forhold mellem Styrke og Indhold. En Følelses
Indhold er de Forestillinger, til hvilke den er knyttet, og som giver
den dens ejendommelige Beskaffenhed. Glæde og Sorg nuanceres og
varieres efter det Forestillingsindhold, der betinger dem. Nu er det
klart, at med det snævre Omfang følger ogsaa let et snævert Indhold.
I de smaa Kredse antager Følelseslivet let en rent elementær Ka-
rakter, kommer til at mangle bestemt og mangesidig Udformning.
Ved Isolation fra Livet i de større Kredse indtræder da --- selv om
Følelsens Styrke bevares, hvilket paa Grund af Vanens og Ensformig-
hedens Indflydelse ikke altid sker --- let en aandelig Fattigdom.
Følelsen bliver da blind. Det Indhold, den har, og som maaske fast-
holdes med stor Inderlighed og Lidenskab, stammer væsenligt fra
Fortiden. Familiens Husguder ere de forbigangne Tiders Heroer;
men det er ikke sagt, at de egne sig til at være Førere i de nye
Tider. Hjemmets konservative Karakter betinges herved. Inderlig-
heden, Trygheden og Hvilen bliver da til Træghed og afføder Uvillie
og Frygt over for det nye. Hvad der ellers udgør Hjemmets Styrke,
bliver her dets Svaghed. Alle Tiders Historie frembyder Erfaringer
i denne Retning. Her ville vi kun fremdrage et Par særligt iøjne-
faldende Træk. --- Nyligt har en japanesisk Forfatter, Tokiwo Yokoi,
i Anledning af Krigen mellem Kina og Japan anstillet en Sammen-
ligning mellem de to Folks herskende Moral og har herved lagt stor
Vægt paa, at den kinesiske Moral væsenligt er en Familiemoral:
»Det var den sønlige Pietets Princip, som var Grundlaget for Kine-
sernes Moral, snarere end Principet om Loyalitet over for Fyrsten.«
Denne Retning slog den kinesiske Etik saa meget mere ind paa, som
Dynastierne i Kina have skiftet meget hyppigt. Familieforholdet blev
da det eneste Baand, der kunde hævde Enheden og Kontinuiteten i
Samfundet, og det med saadan Fasthed, at »disse tre Hundrede Mil-
lioner Mennesker have bevaret deres Traditioner og Skikke ufor-
andrede i tre Tusinde Aar og trods tre og tyve Dynastiforandringer.«
Men de Dyder, der udmærke Kineserne, ere derfor ogsaa kun det
private og huslige Livs Dyder, hvorimod det skorter paa Patriotisme,
paa Sans for og Ærlighed i offentlige Forhold. Deraf forklares Van-
skeligheden ved Optagelse og Gennemførelse af Reformer i Kina,
medens Japanerne, der gennem lange Tider vare øvede i at under-
ordne private og huslige Interesser under Statens Krav, med største
Iver og Opofrelse slog ind paa Fremskridtets Vej. --- Kyndige Iagt-
tagere af Forholdene i Tyrkiet og i Persien have fundet, at de orien-
talske Haremer maaske yde den vigtigste Modstand mod en ny Kulturs
Indtrængen. »Medens der,« siger Vambéry (\emph{Der Islam des
neunzehnten Jahrhunderts}), »i Selamlik (den Del af Huset, der er
tilgængelig for Mændene og den fremmede Verden) hist og her viser
sig Spor af Fornyelse og \emph{à la franca} saa at sige er blevet et
Løsen, er i Haremet, den muhamedanske Families Helligdom, alt
blevet ved det gamle, og der er ikke sket den mindste Forandring i
Møblementet, Sæd, Skik og Husorden, saa lidt som i Dragt, Talebrug
og Tænkemaade\ldots\ Hos nogle forstokkede Ortodokse er Haremet i
Virkeligheden blevet et lykkeligt Tilflugtssted, medens det selvfølgelig
er en vældig Hæmsko for de Forkæmpere for det nye, der mene det
alvorligt. Det nye, der indføres i Selamlik, bliver i Haremet enten
ignoreret eller revet ned igen af ubændig Fanatisme.« --- Det mangler
ikke i Europa paa Tendenser i lignende Retning. Man har saaledes
--- baade i protestantiske og i katolske Lande --- opstillet den
Paastand, at Skolen skal være afhængig af Hjemmet. »Hjem, Stat
og Kirke«, sagde Biskop Ketteler, »maa have Skoler, der svare til
deres Aand og deres Fordringer.« Der mangler her den supplerende
Sætning, at Skolen --- i videste Forstand: Samfundet til Fremme og
Tilegnelse af Videnskaben --- atter skal virke tilbage paa den Aand
og de Fordringer, der raade i Hjem, Stat og Kirke, --- at der er en
større Horisont, som Hjemmets (og de andre nævnte Samfunds) Tendens
til Afslutning let fører til at begrænse. Hjemmets ofte lumre og for-
gemte Luft trænger til at renses gennem friske Strømninger. I saa-
danne Tilfælde kæmper man nu egentligt ikke for Hjemmet, men
man bruger Hjemmet som Middel til at bevare Dogmer, som ikke
paa anden Maade kunne holdes oppe. Men det Kulturindhold, paa
hvilket Hjemmet skal tære, vil ved denne Fremgangsmaade blive
ringere og vil stivne i døde Former.

Det vil i Længden ikke blive muligt at bevare Styrken og Inder-
ligheden, naar Indholdet paa denne Maade skrumper sammen. Gen-
tagelsen vil øve sin sløvende Magt, og aandelig Død blive den sidste
Virkning. Al levende Følelse forudsætter til sin Bestaaen Rytme og
Modsætning; hvor der slet intet nyt er at opdage, gaar Følelsen over
til passiv Vane. Derfor vil Hjemmets Karakter og Værdi bero paa
Rigdommen i de Menneskers Natur, der udgøre Hjemmet. Jo ejen-
dommeligere og flersidigere de hver for sig ere, des gunstigere Be-
tingelser vil der være for nye Opdagelser inden for den snævre Kreds,
og des uafhængigere vil Hjemmet kunne blive af den store Verden,
uden at miste Livet og Friskheden.

Men det er nu netop ikke altid, at den enkelte Personlighed i
Hjemmet finder de rette Betingelser for sin Udvikling og møder den
Forstaae\-lse, uden hvilken dens Ejendommelighed ikke kan udfolde
sig. I Hjemmet bliver det naturligt det fælles, der raader. Natur-
grunden, som bærer Familien, holder de enkelte Individualiteter fast
i deres endnu ubestemte Form; i den bestemtere Udformning af den
enkelte Individualitet ligger jo en Mulighed for, at den umiddelbare
Sammenhæng og den instinktmæssige gensidige Forstaae\-lse skal op-
høre. Der opstaar naturligt en Sky for at bryde Fælleslivet, for at
bøje af fra de vante Spor, ad hvilke det tilvante, paa uskrevne Love
byggede Livssyn fører Familiens Medlemmer, og den enkelte hæmmes
let derved i sin selvstændige Udvikling. I sin mest prosaiske Form
viser denne Sky sig som Generthed. En ejendommelig Art af
Undseelse, der træder i Kraft over for de nærmeste, men som kan
falde bort over for fjernere staaende. Dette Forhold danner en
Analogi til den Kendsgerning, at Elskovsfølelsen paa højere Trin af
Menneskeliv ikke vaagner inden for Søskendeforholdet, men kræver
en ny og fremmed Genstand for at vækkes. Kun over for det nye
og fremmede kommer Individualiteten ret frem. Dertil kommer, at
inden for Hjemmets Kreds kunne baade de uheldige og de lykkelige
Elementer i Karakteren trives. Hvad der let afslibes under Berøring
med det mere skaanselløse fremmede, kan under Hjemmets lune og
bløde Forhold brede sig og drage Opmærksomheden bort fra det
væsenlige og værdifulde. I hvert Tilfælde træder inden for Hjemmet
det ringe og det store ofte frem i uklare Blandingsforhold, der gør
baade Overvurdering og Undervurdering mulig. Snart bliver den
enkelte uden Grund anset som Profet af sine nærmeste, snart ikke
agtet som Profet, skønt han er det.

\section*{IV.}

\noindent Hjemmets Fremtid vil bero paa, om de Farer, som her ere paa-
pegede, kunne undgaas. Skal Hjemmet hævde sin Værdi for det
menneskelige Liv, maa det træde i et aabent og positivt Forhold til
de større Samfund og de Interesser, der røre sig i disse, og det maa
tillige inden for sin egen Kreds tilvejebringe saa gunstige Betingelser
som muligt for sine enkelte Medlemmers ejendommelige Udvikling.
Fra en dobbelt Side kan Hjemmet angribes: fra socialistisk Side
som isolerende, --- fra anarkistisk Side som hæmmende. Den
menneskelige Udvikling gaar paa en Gang i Retning af større Hori-
sonter, mere universelle Virkekredse, og i Retning af bestemtere Ud-
formning af de individuelle Ejendommeligheder.

Men der er ingen Grund til at tro, at Hjemmet ikke skulde kunne
løse de Opgaver, der stilles til det. Det har hidtil fulgt Mennesket
paa den lange Vej fra dyrisk Raahed til socialt og aandeligt Men-
neskeliv, og selv om Husguderne skifte, og maaske for Fremtiden
ville skifte hyppigere end før, saa vil deres Plads ikke komme til at
staa tom.

Det er en Naturnødvendighed, at vi leve i en lille Verden, før den
store kan blive til for os. Og jo mere Hjemmet virkeligt selv er en
lille Verden, des mere kan Livet der blive en Forskole for Livet i
den store Verden, --- samtidigt med, at det bevarer sin selvstændige
Værdi som Inderlighedens og Tryghedens Sted. Der er noget, som
er fælles for enhver Verden, den være lille eller stor. En almindelig
Lov raader og bestemmer de enkelte Elementers Sted og Virken; det
enkelte Element maa underordne sig det heles Orden, men besjæles
tillige af de forbindende Kræfter, der holde det hele sammen. Inden
for Familien lever den enkelte sig uvilkaarligt ind i den Forestilling,
at han er en blandt mange; han lærer Hengivelsen, Resignationen
og Retfærdigheden at kende, og uden Støj, uden at nogen pinlig
Modstand behøver at føles, hindres hans Egoismes ubændige Udfoldelse.
De første Genstande for Frygt og Kærlighed, for Beundring og Tak-
nemmelighed, for Agtelse og Ærefrygt finder han inden for Hjemmets Kreds.

Allerede inden for den lille Verden --- naar denne ikke isolerer
sig --- vil den enkelte gøre Erfaringer om større Forholds Betydning.
Bølgegangen fra Verdenshavet fortsættes ofte ind i Hjemmets stille
Vig. Nationale Ulykker f.\ eks.\ berøre tidligt Barnet gennem den Sorg,
de vække hos de ældre; til sin Forundring ser det, at hvad der sker
i det fjerne, betager de nærmeste, som var det deres eget Ve og
Vel, det drejede sig om. Den første Følelse af at høre til et større
Hele end det, der udgøres af Forældre og Søskende, Slægt og Venner,
grundlægges maaske i et saadant Øjeblik. Gennem Forholdet til de
første Genstande for dets Kærlighed og Beundring lever Barnet sig
uvilkaarligt ind i de Interesser, for hvilke disse dets Forbilleder alle-
rede leve. Men de ældre staa ikke blot som Mellemled mellem
Barnet og de større Interessekredse. De staa ogsaa som Mellemled
mellem Fortiden og Fremtiden. Deres Opgave er ikke blot at over-
levere Fortidens Værdier, men ogsaa at forberede Fremtidens.
Hjemmet har her den store Betydning at slaa Bro mellem de for-
skellige Generationer og gøre en Forstaae\-lse mellem dem mulig, der
vilde udelukkes, dersom den nye Slægt ikke voksede op inden for
den ældre og under Indflydelse af den mest umiddelbare og mest
naturbestemte Sympati, som i det hele kan forbinde Mennesker.
Striden mellem det gamle og det nye vilde være skarpere og tillige
mindre frugtbar, hvis den ikke i dæmpet Form forberedtes inden for
Hjemmets Vægge.

Der er tillige inden for Hjemmet de gunstigste Betingelser for
personlig Udvikling, især i den Periode, hvor Formerne endnu ere
ubestemte, og hvor mange Skud skyde frem, som vel ikke ere leve-
dygtige, men som dog ikke voldsomt kunne afskæres uden Fare for
Livet. Der ligger fremdeles i Sympatiens Natur, saaledes som den
kan fremtræde inden for Hjemmet, en Interesse og en Evne til For-
staae\-lse for rent individuelle Rørelser, som kun under uheldige
Forhold udelukkes, naar Blindhed og Begrænsning kommer til at
raade. I hvert Tilfælde vil individuel Ejendommelighed næppe finde
saa gode Livsbetingelser som her. Der vil ogsaa i denne Henseende
stilles større Fordringer til Fremtidens Hjem. Det er jo nu --- for
at fremhæve et enkelt vigtigt Punkt --- ikke blot Manden, hvis Per-
sonlighed kræver frit Spillerum for individuelle Ejendommeligheder.
Ogsaa for Kvindens Vedkommende rejses nu den samme Fordring.
Ikke blot ønske Kvinderne nu, som det ogsaa i tidligere Tider, f.\ eks.\
i Renaissanceperioden, kunde ske, Del i den almindelige Dannelse
og i Tidens Tankegang, men ogsaa Mulighed for anden Kaldsvirk-
somhed end den, der øves paa selve det huslige Omraade. Arbejdsdelin-
gens store Lov vil ogsaa ytre sig for den kvindelige Dannelses
og det kvindelige Arbejdes Vedkommende. Medens Kvinden inden
for Hjemmet --- baade som Moder, Hustru og Søster --- kunde siges
at repræsentere det udelt menneskelige, de blivende Livskilder for
Inderlighed og Kærlighed, saa truer nu Arbejdsdelingens og Ensidig-
hedens Lov ogsaa dem. Mange se heri en stor Fare for Hjemmets
Fremtid. Men hvis Hjemmet vil kunne løse den Opgave at lade den
personlige Ejendommelighed røre sig friere end hidtil for Mandens
Vedkommende, vil den tilsvarende Opgave ogsaa nok kunne løses for
Kvindens Vedkommende. Hjemmet har saa dyb og fast en Natur-
grund, at det sikkert vil bestaa, selv om den hidtilværende Arbejds-
deling mellem de to Køn afløses af en anden Ordning. Naturen har
her ikke sagt sit sidste Ord.

Som alt andet i Verden maa ogsaa Hjemmet kæmpe for Livet.
Under saare forskellige Former har det hidtil hævdet sin Betydning,
og det vil ogsaa i Fremtiden kunne godtgøre sig i sin store Betydning:
paa en Gang at være en Inderlighedens og en Frihedens Verden i
langt højere Grad end nogen anden Art af Samfund formaar det.

\end{document}
%%% Local Variables:
%%% mode: latex
%%% TeX-master: t
%%% End:
